\section{Bivariables}
El análisis bivariado permite estudiar la relación entre dos variables de manera conjunta, a
diferencia del análisis univariado que examina cada variable por separado. Este enfoque resulta
útil para identificar si existe algún tipo de asociación, dependencia o patrón entre ambas
variables, y para describir cómo se comportan en conjunto dentro de un mismo conjunto de datos.

\subsection{Cualitativa-Cualitativa}
El más sencillo, comparando dos variables cualitativas. En este caso utilizaremos
\Variable{'part\_time\_job'} y \Variable{'absence\_days'}.


\subsubsection{Tabla de contingencia: part\_time\_job vs absence\_days}
\begin{DatosTabla}{L{0.18\linewidth} C{0.15\linewidth} C{0.17\linewidth} C{0.16\linewidth} C{0.12\linewidth}}{\textbf{¿Trabaja?} & \textbf{No faltó} & \textbf{Faltó una vez} & \textbf{Faltó 2 o más} & \textbf{Total}}[Tabla de contingencia: part\_time\_job vs absence\_days]

Trabaja & 16 & 18 & 24 & 58 \\ \hline
No trabaja & 20 & 13 & 20 & 53 \\ \hline
\textbf{Total} & \textbf{36} & \textbf{31} & \textbf{44} & \textbf{111} \\


\end{DatosTabla}

.

% =====================================================================
%  figures/bivariable-cuali-cuali.tex
%  Barras agrupadas: 'part_time_job' x 'absence_days'.
%  (Reemplaza a img/bivariable_cuali_cuali.png)
% =====================================================================
\begin{figure}[H]
\centering
\begin{tikzpicture}
\begin{axis}[
  ybar, bar width=16pt,
  title={Días de ausencia según trabajo a tiempo parcial},
  xlabel={Situación de faltas},
  ylabel={Cantidad de estudiantes},
  ymin=0, ymax=28,
  symbolic x coords={No faltó, Faltó una vez, Faltó 2 o más},
  xtick=data,
  legend style={at={(0.5,-0.18)}, anchor=north, legend columns=-1},
  enlarge x limits=0.3,
  nodes near coords,
]
\addplot+[fill=azul, draw=azul] coordinates {(No faltó,16) (Faltó una vez,18) (Faltó 2 o más,24)};
\addplot+[fill=naranja, draw=naranja] coordinates {(No faltó,20) (Faltó una vez,13) (Faltó 2 o más,20)};
\legend{Trabaja, No trabaja}
\end{axis}
\end{tikzpicture}
\caption{Relación entre trabajo a tiempo parcial y días de ausencia.}
\label{fig:bivariable_cuali_cuali}
\end{figure}


\paragraph{Conclusión:} A partir de la tabla de contingencia no se observa una relación clara entre trabajar y la cantidad de inasistencias. Si bien entre los estudiantes que trabajan hay una mayor cantidad de casos con dos o más faltas, una tendencia similar también se presenta entre quienes no trabajan. En consecuencia, con estos datos no puede afirmarse que tener trabajo impacte de manera significativa en la regularidad de asistencia de los estudiantes.

%---------------------------------------------------------------------------------------------------

\subsection{Cuantitativa-Cuantitativa}
Este tipo de análisis se utiliza cuando ambas variables son cuantitativas, lo que permite evaluar la relación entre ellas mediante medidas estadísticas como el coeficiente de correlación.
\Variable{'weekly\_self\_study\_hours'} y \Variable{'math\_score'}.

\paragraph{Análisis descriptivo de las variables}.

\textbf{Variable independiente $(x)$:} \Variable{weekly\_self\_study\_hours} (Horas de estudio semanal). Es una variable cuantitativa continua.

\textbf{Variable dependiente $(y)$:} \Variable{math\_score} (Puntaje en matemáticas). Es una variable cuantitativa que representa el resultado o salida del estudio.

\textbf{Tamaño de la muestra $(n)$:} 111

\textbf{Horas de estudio $(x)$}
\begin{itemize}
\item Media aritmética $(\bar{x})$: 7,09 horas. Indica que, en promedio, los estudiantes dedican poco más de 7 horas semanales al estudio autónomo.
\item Mediana $(Me)$: 7 horas. El 50\% de los estudiantes estudia 7 horas o menos.
\item Rango: 5 horas. La diferencia entre el máximo (9) y el mínimo (4).
\item Variabilidad: El conjunto es homogéneo $(CV<20\%)$, lo que significa que la media es altamente representativa del grupo.
\end{itemize}

\textbf{Puntaje en matemáticas $(y)$}
\begin{itemize}
\item Media aritmética $(\bar{y})$: 84,57 puntos. El rendimiento promedio del grupo se sitúa cerca de los 85 puntos.
\item Mediana $(Me)$: 85 puntos. La mitad de los estudiantes obtuvo un puntaje de 85 o superior.
\item Moda $(Mo)$: 85 puntos. Es el puntaje más frecuente en la muestra.
\item Rango: 17 puntos. Los puntajes oscilan entre un mínimo de 75 y un máximo de 92.
\end{itemize}
\newpage
\subsubsection{Tabla de contingencia: weekly\_self\_study\_hours vs math\_score}
\begin{DatosTabla}{L{0.20\linewidth} C{0.15\linewidth} C{0.15\linewidth} C{0.15\linewidth} C{0.13\linewidth}}{\textbf{Horas de estudio} & \textbf{75-80} & \textbf{81-85} & \textbf{86-92} & \textbf{Total}}[Tabla de contingencia: weekly\_self\_study\_hours vs math\_score]

4-5 & 7 & 12 & 1 & 20 \\ \hline
6-7 & 3 & 30 & 16 & 49 \\ \hline
8-9 & 1 & 22 & 19 & 42 \\ \hline
\textbf{Total} & \textbf{11} & \textbf{64} & \textbf{36} & \textbf{111} \\


\end{DatosTabla}



% =====================================================================
%  figures/bivariable-cuanti-cuanti.tex
%  Dispersión: 'weekly_self_study_hours' x 'math_score'.
%  (Reemplaza a img/bivariable_cuanti_cuanti.png)
%
%  NOTA: el repositorio no conserva los 111 datos crudos, solo la tabla
%  de contingencia agrupada (horas x puntaje). Esta nube se reconstruyó
%  respetando EXACTAMENTE esas frecuencias por celda, por lo que muestra
%  la misma relación (tendencia positiva) que el gráfico original.
% =====================================================================
\begin{figure}[H]
\centering
\begin{tikzpicture}
\begin{axis}[
  dispersion,
  title={Relación entre horas de estudio semanal y puntaje en matemáticas},
  xlabel={Horas de estudio (\Variable{weekly\_self\_study\_hours})},
  ylabel={Puntaje (\Variable{math\_score})},
  xmin=3, xmax=10, ymin=70, ymax=100,
  legend style={at={(0.98,0.02)}, anchor=south east},
]
\addplot+[only marks, mark=*, mark size=1.7pt, azul] coordinates {
  % Horas 4-5
  (4,76) (5,78) (4,79) (5,76) (4,77) (5,79) (4,78)
  (4,82) (5,84) (4,83) (5,82) (4,84) (5,83) (4,81) (5,85) (4,85) (5,81) (4,82) (5,84)
  (5,92)
  % Horas 6-7
  (6,75) (7,78) (6,79)
  (6,82) (7,84) (6,83) (7,82) (6,84) (7,83) (6,81) (7,85) (6,85) (7,81)
  (6,82) (7,84) (6,83) (7,82) (6,84) (7,83) (6,81) (7,85) (6,82) (7,84)
  (6,83) (7,82) (6,84) (7,83) (6,81) (7,85) (6,82) (7,84) (6,83) (7,82)
  (6,88) (7,89) (6,87) (7,88) (6,89) (7,87) (6,90) (7,86) (6,86) (7,90)
  (6,88) (7,89) (6,87) (7,88) (6,89) (7,92)
  % Horas 8-9
  (8,78)
  (8,82) (9,84) (8,83) (9,82) (8,84) (9,83) (8,81) (9,85) (8,85) (9,81)
  (8,82) (9,84) (8,83) (9,82) (8,84) (9,83) (8,81) (9,85) (8,82) (9,84) (8,83) (9,82)
  (8,88) (9,89) (8,87) (9,88) (8,89) (9,87) (8,90) (9,86) (8,86) (9,90)
  (8,88) (9,89) (8,87) (9,88) (8,89) (9,87) (8,90) (9,86) (8,88)
};
\end{axis}
\end{tikzpicture}
\caption{Relación entre horas de estudio semanal y puntaje en matemáticas.}
\label{fig:bivariable_cuanti_cuanti}
\end{figure}


\paragraph{Análisis bivariado (relación entre variables)}.

\textbf{Patrón de relación:}Al observar los datos agrupados, se identifica una tendencia positiva. A medida que aumentan las horas de estudio semanales (pasando del bloque de 4-5 horas al de 8-9 horas), los puntajes de matemáticas tienden a concentrarse de forma notable en los niveles más altos de la muestra (entre 81 y 92 puntos).

\textbf{Dispersión en la relación:} Se observa que estudiantes con una dedicación horaria baja (por ejemplo, 5 horas) pueden alcanzar puntajes máximos de la muestra (92 puntos). Esto indica que, si bien existe una clara relación lineal positiva, operan otros factores asociados que influyen en el rendimiento académico además del tiempo de estudio autónomo.
%---------------------------------------------------------------------------------------------------
\subsection{Cuantitativa-Cualitativa}
Este tipo de análisis se utiliza cuando una variable es cuantitativa y la otra es cualitativa, lo que permite evaluar cómo varía la variable cuantitativa en función de las categorías de la variable cualitativa. En este caso, se analizará la relación entre \Variable{'gender'} y \Variable{'biology\_score'}.

\subsubsection{Medidas resumen de biology\_score por gender}
% Nota: tabla ancha (12 columnas) heredada del resumen de SPSS. Se reduce
% la tipografía y el espaciado para que entre en el ancho de página.
\footnotesize
\setlength{\tabcolsep}{3pt}
\begin{DatosTabla}{l l c c c c c c c c c c}{\textbf{Gender} & \textbf{Variable} & \textbf{n} & \textbf{Media} & \textbf{D.E.} & \textbf{Var (n-1)} & \textbf{CV} & \textbf{Mín} & \textbf{Máx} & \textbf{Mediana} & \textbf{Q1} & \textbf{Q3}}[Medidas resumen de biology\_score por gender]

female & Biology\_Score & 58 & 83,62 & 2,61 & 6,80 & 3,12 & 79,00 & 90,00 & 83,50 & 82,00 & 86,00 \\ \hline
male & Biology\_Score & 53 & 84,17 & 2,26 & 5,11 & 2,68 & 79,00 & 90,00 & 85,00 & 83,00 & 86,00 \\


\end{DatosTabla}
\normalsize





\subsubsection{Comparación de métricas muestrales por género}
\begin{table}[H]
\centering
\begin{tabular}{lcc}
\toprule
\textbf{Métrica} & \textbf{Femenino (female)} & \textbf{Masculino (male)} \\
\midrule
Tamaño de submuestra ($n$) & 58 estudiantes & 53 estudiantes \\
Media aritmética ($\bar{x}$) & 83.62 puntos & 84.17 puntos \\
Desvío estándar ($S$) & 2.61 puntos & 2.26 puntos \\
Coeficiente de variación ($cv\%$) & 3.12\% & 2.68\% \\
\bottomrule
\end{tabular}
\end{table}

\paragraph{Interpretación:} La distribución del grupo es equilibrada, ya que ambos géneros presentan tamaños de submuestra similares. El rendimiento promedio es muy parecido entre mujeres y hombres, con una diferencia menor a 0.6 puntos. Además, la dispersión de las notas es baja en ambos casos, y los coeficientes de variación indican que los grupos son altamente homogéneos ($cv\% < 20\%$).

% =====================================================================
%  figures/bivariable-cuali-cuanti.tex
%  Boxplot: 'gender' x 'biology_score'.
%  (Reemplaza a img/bivariable_cuali_cuanti.png)
%
%  Valores desde la tabla de SPSS (tables/bivariable-genero-biologia):
%    female: min 79, Q1 82, mediana 83,5, Q3 86, max 90
%    male:   min 79, Q1 83, mediana 85,   Q3 86, max 90
% =====================================================================
\begin{figure}[H]
\centering
\begin{tikzpicture}
\begin{axis}[
  boxplot/draw direction=y,
  title={Distribución de \Variable{biology\_score} según género},
  ylabel={Puntaje de biología},
  ymin=76, ymax=92,
  xtick={1,2},
  xticklabels={female, male},
  xlabel={Gender},
  grid=major,
]
  % female
  \addplot+[boxplot prepared={
      lower whisker=79, lower quartile=82, median=83.5,
      upper quartile=86, upper whisker=90},
      fill=azulclaro, draw=azul] coordinates {};
  % male
  \addplot+[boxplot prepared={
      lower whisker=79, lower quartile=83, median=85,
      upper quartile=86, upper whisker=90},
      fill=naranja!40, draw=naranja] coordinates {};
\end{axis}
\end{tikzpicture}
\caption{Relación entre género y puntaje en biología.}
\label{fig:bivariable_cuali_cuanti}
\end{figure}


\paragraph{Detalle para la imagen}
\begin{itemize}
    \item \textbf{Mujeres (Female):} la media se ubica exactamente en 83,62 puntos, lo que indica un rendimiento promedio alto y muy similar al de los hombres.
    \item \textbf{Hombres (Male):} la media se encuentra en 84,17 puntos, apenas 0,55 puntos por encima del valor de las mujeres, lo que sugiere una diferencia mínima.
    \item \textbf{Mediana:} en el caso de las mujeres, la línea se sitúa ligeramente por debajo del punto central, cerca de 83,5 puntos; en el caso de los hombres, la línea se ubica por encima del punto, en 85,5 puntos, indicando que más de la mitad de los varones superó esa nota.
\end{itemize}
